\makeatletter
\def\input@path{{./}{../../../preamble/}{book2026-2027/preamble/}}
\makeatother
\input{preamble.tex}
\usepackage{needspace}
\hypersetup{pdftitle={Grade 10 Term 2 Discrete Probability}}
\setlength{\emergencystretch}{2em}
\title{Grade 10, Term 2, Lesson 1, Criterion C4\\Binomial Counts in Fixed Financial Trials}
\author{}
\date{}
\begin{document}
\maketitle
\setcounter{tocdepth}{1}
\setcounter{secnumdepth}{2}
\phantomsection\label{toc}
\tableofcontents
\section*{Evaluated criterion}
\textbf{C4:} Applies the binomial distribution to fixed-trial financial situations by verifying its conditions and calculating probabilities using the binomial probability formula.
\section{Introduction}\label{sec:t2-c4-introduction}
All data in this activity are hypothetical teaching models. Probabilities have no monetary units. Keep exact expressions where useful and round final probabilities to four decimal places.
\subsection{A fixed set of opportunities}
A school shop sends four independent offers. Each visitor either purchases (success) or does not, with constant purchase probability $p$. Let $X$ count purchases. The number of trials $n=4$ is fixed before results are observed. The support is $\{0,1,2,3,4\}$ purchases. Probability has no monetary units; revenue might be a separate variable $Y=5X$ USD.
\textbf{Check.} Does the shop stop after its first sale? \emph{Suggested answer:} No. A binomial count observes all four fixed trials.


\subsection{Four conditions}
Use binomial only with a fixed number $n$ of binary trials, independent results and constant success probability $p$. Independent means one result does not change the probability of another. A common economic shock may break independence. Changing customer groups may break constant probability.
\textbf{Check.} Are default events automatically independent? \emph{Suggested answer:} No. A shared economic shock can link them.


\subsection{Count the arrangements}
Exactly two successes in three trials can occur as SSF, SFS or FSS. Each sequence has probability $p^2(1-p)$. There are three such arrangements, so add their probabilities. In general their number is the combination $\binom nk$, read as n choose k:
\[\binom nk=\frac{n!}{k!(n-k)!}.\]
The factorial $n!$ means multiply positive integers from 1 to $n$: $4!=4\cdot3\cdot2\cdot1=24$, and $0!=1$. Thus $\binom42=24/(2\cdot2)=6$. Order matters for a sequence, but the combination counts all sequences with the same number of successes.
\textbf{Check.} Why is $p^2(1-p)$ alone insufficient for two successes in three trials? \emph{Suggested answer:} It describes only one arrangement; there are three.

\textbf{Check.} What does $\sum_{j=0}^{2}P(X=j)$ mean? \emph{Suggested answer:} Add $P(X=0)+P(X=1)+P(X=2)$; it is an at-most probability.

\section{General worked problems}
\Needspace{14\baselineskip}
\subsection{Introductory general problem: exactly two successes}
Three independent binary trials each have success probability $1/2$.

\textbf{Questions/Tasks}
\begin{enumerate}
\item Calculate $P(X=2)$, defining success as the trial gives the designated S outcome. 
\end{enumerate}
\subsubsection*{Complete worked solution}
\begin{center}
\begin{tabular}{llp{0.48\linewidth}}
\toprule
Symbol & Value & Meaning in this problem \\ \midrule
$n$ & $3$ trials & Fixed number of trials \\
$p$ & $0.5$ & Success probability per trial \\
$k$ & $2$ successes & Requested success count \\
\bottomrule
\end{tabular}
\end{center}

\par\noindent\begin{minipage}{\linewidth}
\textbf{Step 1: Define the variable.} $X$ counts successes among 3 trials. A success means the trial gives the designated S outcome. Its support is $\{0,1,\ldots,3\}$ successes; money is a separate consequence.
\end{minipage}\par\medskip
\par\noindent\begin{minipage}{\linewidth}
\textbf{Step 2: State parameters and check conditions.} $n=3$ is fixed, $p=0.5$ and $1-p=0.5$. Under the stated independence, binary-outcome and constant-probability assumptions, $X\sim\operatorname{Bin}(3,0.5)$.
\end{minipage}\par\medskip
\par\noindent\begin{minipage}{\linewidth}
\textbf{Step 3: Translate and calculate.} The requested event is $ X=2$. First use
\[P(X=k)=\binom nk p^k(1-p)^{n-k}.\]
Then substitute:
\[P(X=2)=\binom{3}{2}(0.5)^2(0.5)^{1}\approx 0.3750.\]
\end{minipage}\par\medskip
\par\noindent\begin{minipage}{\linewidth}
\textbf{Step 4: Check and interpret.} The answer is between 0 and 1; the full binomial PMF totals $(p+(1-p))^n=1$. Under this hypothetical model, the probability of the requested success count is 0.3750. 
\end{minipage}\par\medskip
\subsection{Bridge to the general model}
\[
X\sim\operatorname{Bin}(n,p),\qquad P(X=k)=\binom nk p^k(1-p)^{n-k},\quad k=0,\ldots,n.
\]
The powers describe one sequence: $k$ successes and $n-k$ failures. The combination counts its possible arrangements. For at most $b$, add terms from 0 to $b$. For at least $a$, use $1-P(X\leq a-1)$; the complement ends at $a-1$, not $a$. For between two values, state whether both endpoints are included.
\Needspace{14\baselineskip}
\subsection{Advanced general problem: a central interval}
Six independent binary trials have constant success probability 0.4.

\textbf{Questions/Tasks}
\begin{enumerate}
\item Calculate $P(2\leq X\leq4)$, defining success as the designated event occurs. 
\end{enumerate}
\subsubsection*{Complete worked solution}
\par\noindent\begin{minipage}{\linewidth}
\textbf{Step 1: Define the variable.} $X$ counts successes among 6 trials. A success means the designated event occurs. Its support is $\{0,1,\ldots,6\}$ successes; money is a separate consequence.
\end{minipage}\par\medskip
\par\noindent\begin{minipage}{\linewidth}
\textbf{Step 2: State parameters and check conditions.} $n=6$ is fixed, $p=0.4$ and $1-p=0.6$. Under the stated independence, binary-outcome and constant-probability assumptions, $X\sim\operatorname{Bin}(6,0.4)$.
\end{minipage}\par\medskip
\par\noindent\begin{minipage}{\linewidth}
\textbf{Step 3: Translate and calculate.} The requested event is $ 2\leq X\leq4$. First use
\[P(a\leq X\leq b)=\sum_{j=a}^{b}\binom nj p^j(1-p)^{n-j}.\]
Then substitute:
\[P(2\leq X\leq4)=\sum_{j=2}^{4}\binom{6}j(0.4)^j(0.6)^{6-j}\approx 0.7258.\]
\end{minipage}\par\medskip
\par\noindent\begin{minipage}{\linewidth}
\textbf{Step 4: Check and interpret.} The answer is between 0 and 1; the full binomial PMF totals $(p+(1-p))^n=1$. Under this hypothetical model, the probability of the requested success count is 0.7258. 
\end{minipage}\par\medskip
\section{Notation reference and practice procedure}
\begin{center}
\begin{tabular}{lp{0.72\linewidth}}
\toprule
Symbol & Meaning and units \\ \midrule
$n,k$ & Fixed trial count and realised success count \\
$p,1-p$ & Success and failure probabilities \\
$\binom nk$ & Number of success-position arrangements \\
$n!$ & Factorial, with $0!=1$ \\
$\sum_{j=a}^b$ & Add one term for every integer $j$ from $a$ to $b$ \\
\bottomrule
\end{tabular}
\end{center}
\begin{enumerate}
\item Define the random variable in words and give its units and support.
\item List parameters and check that the model assumptions match the situation.
\item Translate ordinary language into an event, paying attention to endpoints.
\item Write the symbolic formula first, then substitute without rounding intermediate values.
\item Check bounds, total probability or a complementary calculation; interpret the answer in context.
\end{enumerate}
\Needspace{19\baselineskip}
\section{Foundational contextual practice}\label{sec:t2-c4-foundational}
\Needspace{14\baselineskip}
\subsection{Loan defaults}
Five independent loans each default within one year with probability 0.1. Default is success mathematically.

\textbf{Questions/Tasks}
\begin{enumerate}
\item Calculate $P(X=1)$, defining success as a loan defaults within the year. 
\end{enumerate}
\subsubsection*{Complete worked solution}
\par\noindent\begin{minipage}{\linewidth}
\textbf{Step 1: Define the variable.} $X$ counts successes among 5 trials. A success means a loan defaults within the year. Its support is $\{0,1,\ldots,5\}$ successes; money is a separate consequence.
\end{minipage}\par\medskip
\par\noindent\begin{minipage}{\linewidth}
\textbf{Step 2: State parameters and check conditions.} $n=5$ is fixed, $p=0.1$ and $1-p=0.9$. Under the stated independence, binary-outcome and constant-probability assumptions, $X\sim\operatorname{Bin}(5,0.1)$.
\end{minipage}\par\medskip
\par\noindent\begin{minipage}{\linewidth}
\textbf{Step 3: Translate and calculate.} The requested event is $ X=1$. First use
\[P(X=k)=\binom nk p^k(1-p)^{n-k}.\]
Then substitute:
\[P(X=1)=\binom{5}{1}(0.1)^1(0.9)^{4}\approx 0.3281.\]
\end{minipage}\par\medskip
\par\noindent\begin{minipage}{\linewidth}
\textbf{Step 4: Check and interpret.} The answer is between 0 and 1; the full binomial PMF totals $(p+(1-p))^n=1$. Under this hypothetical model, the probability of the requested success count is 0.3281. 
\end{minipage}\par\medskip
\secInicio[][sec:t2-c4-foundational]
\Needspace{14\baselineskip}
\subsection{On-time invoices}
Four invoices are paid independently, each with constant probability 0.8 of being on time.

\textbf{Questions/Tasks}
\begin{enumerate}
\item Calculate $P(X=3)$, defining success as an invoice is paid on time. 
\end{enumerate}
\subsubsection*{Complete worked solution}
\par\noindent\begin{minipage}{\linewidth}
\textbf{Step 1: Define the variable.} $X$ counts successes among 4 trials. A success means an invoice is paid on time. Its support is $\{0,1,\ldots,4\}$ successes; money is a separate consequence.
\end{minipage}\par\medskip
\par\noindent\begin{minipage}{\linewidth}
\textbf{Step 2: State parameters and check conditions.} $n=4$ is fixed, $p=0.8$ and $1-p=0.2$. Under the stated independence, binary-outcome and constant-probability assumptions, $X\sim\operatorname{Bin}(4,0.8)$.
\end{minipage}\par\medskip
\par\noindent\begin{minipage}{\linewidth}
\textbf{Step 3: Translate and calculate.} The requested event is $ X=3$. First use
\[P(X=k)=\binom nk p^k(1-p)^{n-k}.\]
Then substitute:
\[P(X=3)=\binom{4}{3}(0.8)^3(0.2)^{1}\approx 0.4096.\]
\end{minipage}\par\medskip
\par\noindent\begin{minipage}{\linewidth}
\textbf{Step 4: Check and interpret.} The answer is between 0 and 1; the full binomial PMF totals $(p+(1-p))^n=1$. Under this hypothetical model, the probability of the requested success count is 0.4096. 
\end{minipage}\par\medskip
\secInicio[][sec:t2-c4-foundational]
\Needspace{19\baselineskip}
\section{Intermediate contextual practice}\label{sec:t2-c4-intermediate}
\Needspace{14\baselineskip}
\subsection{Insurance claims at most one}
Six independent policies each have probability 0.05 of a claim this month.

\textbf{Questions/Tasks}
\begin{enumerate}
\item Calculate $P(X\leq1)$, defining success as a policy has at least one claim this month. 
\end{enumerate}
\subsubsection*{Complete worked solution}
\par\noindent\begin{minipage}{\linewidth}
\textbf{Step 1: Define the variable.} $X$ counts successes among 6 trials. A success means a policy has at least one claim this month. Its support is $\{0,1,\ldots,6\}$ successes; money is a separate consequence.
\end{minipage}\par\medskip
\par\noindent\begin{minipage}{\linewidth}
\textbf{Step 2: State parameters and check conditions.} $n=6$ is fixed, $p=0.05$ and $1-p=0.95$. Under the stated independence, binary-outcome and constant-probability assumptions, $X\sim\operatorname{Bin}(6,0.05)$.
\end{minipage}\par\medskip
\par\noindent\begin{minipage}{\linewidth}
\textbf{Step 3: Translate and calculate.} The requested event is $ X\leq1$. First use
\[P(a\leq X\leq b)=\sum_{j=a}^{b}\binom nj p^j(1-p)^{n-j}.\]
Then substitute:
\[P(X\leq1)=\sum_{j=0}^{1}\binom{6}j(0.05)^j(0.95)^{6-j}\approx 0.9672.\]
\end{minipage}\par\medskip
\par\noindent\begin{minipage}{\linewidth}
\textbf{Step 4: Check and interpret.} The answer is between 0 and 1; the full binomial PMF totals $(p+(1-p))^n=1$. Under this hypothetical model, the probability of the requested success count is 0.9672. 
\end{minipage}\par\medskip
\secInicio[][sec:t2-c4-intermediate]
\Needspace{14\baselineskip}
\subsection{At least one business sale}
Eight independent enquiries each have constant sale probability 0.25.

\textbf{Questions/Tasks}
\begin{enumerate}
\item Calculate $P(X\geq1)$, defining success as an enquiry becomes a sale. 
\end{enumerate}
\subsubsection*{Complete worked solution}
\par\noindent\begin{minipage}{\linewidth}
\textbf{Step 1: Define the variable.} $X$ counts successes among 8 trials. A success means an enquiry becomes a sale. Its support is $\{0,1,\ldots,8\}$ successes; money is a separate consequence.
\end{minipage}\par\medskip
\par\noindent\begin{minipage}{\linewidth}
\textbf{Step 2: State parameters and check conditions.} $n=8$ is fixed, $p=0.25$ and $1-p=0.75$. Under the stated independence, binary-outcome and constant-probability assumptions, $X\sim\operatorname{Bin}(8,0.25)$.
\end{minipage}\par\medskip
\par\noindent\begin{minipage}{\linewidth}
\textbf{Step 3: Translate and calculate.} The requested event is $ X\geq1$. First use
\[P(X\geq k)=1-\sum_{j=0}^{k-1}\binom nj p^j(1-p)^{n-j}.\]
Then substitute:
\[P(X\geq1)=1-\sum_{j=0}^{0}\binom{8}j(0.25)^j(0.75)^{8-j}\approx 0.8999.\]
\end{minipage}\par\medskip
\par\noindent\begin{minipage}{\linewidth}
\textbf{Step 4: Check and interpret.} The answer is between 0 and 1; the full binomial PMF totals $(p+(1-p))^n=1$. Under this hypothetical model, the probability of the requested success count is 0.8999. For at least one, the complementary event is no successes. 
\end{minipage}\par\medskip
\secInicio[][sec:t2-c4-intermediate]
\Needspace{19\baselineskip}
\section{Advanced contextual practice}\label{sec:t2-c4-advanced}
\Needspace{14\baselineskip}
\subsection{Payment fees in a band}
A firm processes ten independent payments, each late with probability 0.2. A late payment costs 5 USD, so total late fees are $Y=5X$ USD.

\textbf{Questions/Tasks}
\begin{enumerate}
\item Calculate $P(2\leq X\leq4)$, defining success as a payment is late. Explain why this is also the event $10\leq Y\leq20$ USD.
\end{enumerate}
\subsubsection*{Complete worked solution}
\par\noindent\begin{minipage}{\linewidth}
\textbf{Step 1: Define the variable.} $X$ counts successes among 10 trials. A success means a payment is late. Its support is $\{0,1,\ldots,10\}$ successes; money is a separate consequence.
\end{minipage}\par\medskip
\par\noindent\begin{minipage}{\linewidth}
\textbf{Step 2: State parameters and check conditions.} $n=10$ is fixed, $p=0.2$ and $1-p=0.8$. Under the stated independence, binary-outcome and constant-probability assumptions, $X\sim\operatorname{Bin}(10,0.2)$.
\end{minipage}\par\medskip
\par\noindent\begin{minipage}{\linewidth}
\textbf{Step 3: Translate and calculate.} The requested event is $ 2\leq X\leq4$. First use
\[P(a\leq X\leq b)=\sum_{j=a}^{b}\binom nj p^j(1-p)^{n-j}.\]
Then substitute:
\[P(2\leq X\leq4)=\sum_{j=2}^{4}\binom{10}j(0.2)^j(0.8)^{10-j}\approx 0.5914.\]
\end{minipage}\par\medskip
\par\noindent\begin{minipage}{\linewidth}
\textbf{Step 4: Check and interpret.} The answer is between 0 and 1; the full binomial PMF totals $(p+(1-p))^n=1$. Under this hypothetical model, the probability of the requested success count is 0.5914. The fee transformation gives $10\leq5X\leq20$ USD, hence $2\leq X\leq4$. Both monetary endpoints are included.
\end{minipage}\par\medskip
\secInicio[][sec:t2-c4-advanced]
\Needspace{14\baselineskip}
\subsection{Supplier defects and a failed assumption}
In version A, seven independently produced items each have defect probability 0.15. In version B, seven items are drawn without replacement from a small lot containing both defective and sound items.

\textbf{Questions/Tasks}
\begin{enumerate}
\item Calculate $P(X\geq2)$, defining success as an item is defective. Calculate for version A; explain why version B does not justify the same binomial model.
\end{enumerate}
\subsubsection*{Complete worked solution}
\par\noindent\begin{minipage}{\linewidth}
\textbf{Step 1: Define the variable.} $X$ counts successes among 7 trials. A success means an item is defective. Its support is $\{0,1,\ldots,7\}$ successes; money is a separate consequence.
\end{minipage}\par\medskip
\par\noindent\begin{minipage}{\linewidth}
\textbf{Step 2: State parameters and check conditions.} $n=7$ is fixed, $p=0.15$ and $1-p=0.85$. Under the stated independence, binary-outcome and constant-probability assumptions, $X\sim\operatorname{Bin}(7,0.15)$.
\end{minipage}\par\medskip
\par\noindent\begin{minipage}{\linewidth}
\textbf{Step 3: Translate and calculate.} The requested event is $ X\geq2$. First use
\[P(X\geq k)=1-\sum_{j=0}^{k-1}\binom nj p^j(1-p)^{n-j}.\]
Then substitute:
\[P(X\geq2)=1-\sum_{j=0}^{1}\binom{7}j(0.15)^j(0.85)^{7-j}\approx 0.2834.\]
\end{minipage}\par\medskip
\par\noindent\begin{minipage}{\linewidth}
\textbf{Step 4: Check and interpret.} The answer is between 0 and 1; the full binomial PMF totals $(p+(1-p))^n=1$. Under this hypothetical model, the probability of the requested success count is 0.2834. In the second version, sampling without replacement from a small lot changes the defect probability after each draw. The binomial assumptions fail; a different sampling model or the full lot composition is needed.
\end{minipage}\par\medskip
\secInicio[][sec:t2-c4-advanced]
\section{Closing reference: key ideas and common errors}
Check fixed $n$, binary outcomes, independence and constant $p$ before calculating. The combination counts arrangements; do not omit it. At least $k$ complements at most $k-1$. Count successes separately from monetary costs. Sampling without replacement and common shocks can invalidate the model.
\secInicio
\end{document}
